\section{Imagens e colunas}

\begin{frame}[fragile]{Inserção de figuras}
\begin{columns}[T]
\begin{column}{0.48\textwidth}
\begin{lstlisting}
\begin{figure}
  \centering
  \includegraphics[
    width=.82\linewidth
  ]{images/NeuralNetwork.png}
  \caption{Representacao de uma rede neural}
  \label{fig:rede-neural}
\end{figure}
\end{lstlisting}
\end{column}
\begin{column}{0.48\textwidth}
\begin{figure}
  \centering
  \includegraphics[width=.82\linewidth]{images/NeuralNetwork.png}
  \caption{Representação de uma rede neural}
  \label{fig:rede-neural}
\end{figure}
\end{column}
\end{columns}
\end{frame}

\begin{frame}[fragile]{Tamanho e proporção da imagem}
\begin{lstlisting}
\includegraphics[width=.7\linewidth]{arquivo.png}
\includegraphics[height=4cm]{arquivo.png}
\includegraphics[width=\linewidth,
                 keepaspectratio]{arquivo.png}
\end{lstlisting}

\begin{itemize}
    \item \verb|width| controla a largura; \verb|height| controla a altura.
    \item \verb|\linewidth| representa a largura disponível no ambiente atual.
    \item \verb|keepaspectratio| evita distorção quando largura e altura aparecem juntas.
\end{itemize}
\end{frame}

\begin{frame}[fragile]{Texto e imagem em colunas}
\begin{columns}[T]
\begin{column}{0.52\textwidth}
\begin{lstlisting}
\begin{columns}[T]
  \begin{column}{.42\textwidth}
    Texto explicativo.
  \end{column}
  \begin{column}{.54\textwidth}
    \includegraphics[width=\linewidth]
      {images/cnn.jpg}
  \end{column}
\end{columns}
\end{lstlisting}
\end{column}
\begin{column}{0.44\textwidth}
As larguras das colunas devem somar menos de \verb|\textwidth| para deixar espaço entre elas.

\vspace{0.7em}
O modificador \verb|[T]| alinha o conteúdo pelo topo.
\end{column}
\end{columns}
\end{frame}

\begin{frame}[plain]
    \centering
    \includegraphics[width=.92\linewidth,height=.88\textheight,keepaspectratio]{images/Pre-processing-visualization.pdf}

    \scriptsize Um slide \texttt{[plain]} oculta título, cabeçalho e rodapé para destacar a figura.
\end{frame}
